
\section{Sound Production in Musical Instruments}

\subsection{String Instruments}

In string instruments (like guitars or violins), sound is produced by vibrating strings.

Standing waves form on the string. The possible wavelengths are determined by the length of the string.

\begin{equation}
    l = \frac{n}{2}\lambda
\end{equation}

\begin{itemize}
    \item $l$: length of the string
    \item $\lambda$: wavelength
\end{itemize}

\paragraph{What affects the pitch?}
\begin{itemize}
    \item Length of the string (shorter $\Rightarrow$ higher pitch)
    \item Tension of the string (higher tension $\Rightarrow$ higher pitch)
    \item Mass of the string (lighter $\Rightarrow$ higher pitch)
\end{itemize}

\subsection{Wind Instruments}

In wind instruments, sound is produced by vibrating air columns.

Depending on whether the ends are open or closed, different standing waves can form.

\paragraph{Important Idea:}
The length of the air column determines which frequencies are possible.

\subsection{Percussion Instruments}

Percussion instruments (like drums) behave differently.

They do not produce simple harmonic frequencies. Instead, their vibrations are more complex.

This is why their sound is often less ``pure'' compared to string or wind instruments.
